In some applications the values that $\sigma_k$ can take in \eqref{eq:switchsys} may depend on $\sigma_{k-1}, \sigma_{k-2},\\\ldots$.
These constraints are often conveniently represented using a \emph{finite automaton} and the JSR under such constraints
is called \emph{constrained joint spectral radius} (CJSR) \cite{dai2012gel}; an example of constrained s\witch{} system is given by \exemref{run1} and its automaton is illustrated by \figref{run}.


\begin{myexem}[Running example]%
  \label{exem:run1}
  We borrow the example of \cite[\Sect4]{philippe2016stability}.
  The set of matrices $\A$ is composed of the following four matrices
  \begin{align*}
    A_1 & = A+B
    \begin{pmatrix}
      k_1 & k_2
    \end{pmatrix},&
    A_2 & = A+B
    \begin{pmatrix}
      0 & k_2
    \end{pmatrix},\\
    A_3 & = A+B
    \begin{pmatrix}
      k_1 & 0
    \end{pmatrix},&
    A_4 & = A.
  \end{align*}
  where $k_1 = -0.49$, $k_2 = 0.27$,
  \[
    A =
    \begin{pmatrix}
      0.94 & 0.56\\
      0.14 & 0.46
    \end{pmatrix} \text{ and }
    B =
    \begin{pmatrix}
      0\\
      1
    \end{pmatrix}.
  \]

  The corresponding automaton is represented by Figure~\ref{fig:run}.

  \begin{figure}[!ht]
    \centering
    \begin{tikzpicture}[x=40,y=18]
       \Vertex{1}
       \NO[unit=4](1){2}
       \SOEA[unit=2](2){3}
       \EA[unit=2](3){4}
       \tikzset{EdgeStyle/.style = {->}}
       \tikzset{LoopStyle/.style = {->}}
       \tikzset{EdgeStyle/.append style = {bend right=18}}
       \Edge[label=1](1)(3)
       \Edge[label=3](1)(2)
       \Edge[label=2](2)(1)
       \Edge[label=1](2)(3)
       \Edge[label=2](3)(1)
       \Edge[label=3](3)(2)
       \Edge[label=4](3)(4)
       \draw[thick,->] (3) to [out=115,in=65,looseness=10] node [midway, fill=white] {1} (3);
       \Edge[label=1](4)(3)
     \end{tikzpicture}
     \caption{Automaton for the running example.
     The numbers on the \arcs{} are their respective labels.
     }
     \label{fig:run}
  \end{figure}
\end{myexem}


The automaton representing the constraints can be represented by a
strongly connected labelled directed graph $\G(\Nodes,\Arcs)$
of \nodes{} $\Nodes$ and \arcs{} $\Arcs$,
possibly with parallel \arcs{}.  The labels are elements of the set
$[m]$ and $\Arcs$ is a subset of $\Nodes \times \Nodes
\times [m]$.  We say that $(u,v,\sigma) \in \Arcs$ if
there is an \arc{} between node $u$ and node $v$ with label $\sigma$.
The iteration~\ref{eq:switchsys} is rewritten as follows to take the automaton into account:
\begin{equation}
  \label{eq:cswitchsys}
  x_k = A_{\sigma_k} x_{k-1}, \qquad (\sigma_1, \ldots, \sigma_k) \text{ are the respective labels of a path in }\G
  .
\end{equation}
